\exercisesection{完备局部环的整闭包}

\begin{enumerate}
\def\labelenumi{\arabic{enumi})}
\tightlist
\item
  设 A 为一个维数为 1 的半局部诺特整环。
\end{enumerate}

\begin{enumerate}
\def\labelenumi{\alph{enumi})}
\item
  A 的整闭包是有限 A-代数，当且仅当 A 的完备化 \(\hat{\mathbf{A}}\) 是约化的。（若 \(\hat{\mathbf{A}}\) 是约化的，使用 § 4 第 2 号定理 2 的推论。反之，将情形化归到 A 是整闭的情形，并仿照 § 4 第 4 号定理 3 的证明进行论证。）
\item
  以下三个条件等价：
\end{enumerate}

\begin{enumerate}
\def\labelenumi{(\roman{enumi})}
\item
  A 是日式环；
\item
  A 是 Nagata 环；
\item
  \(\hat{\mathbf{A}}\) 是约化的，并且若 \(q_{1},\ldots,q_{n}\) 是 \(\hat{\mathbf{A}}\) 的极小素理想，则域 \(\kappa(q_{i})\) 是 A 的分式域的可分扩张。
\end{enumerate}

\begin{enumerate}
\def\labelenumi{\arabic{enumi})}
\setcounter{enumi}{1}
\tightlist
\item
  设 A 为一个维数为 1 的诺特局部环。则以下条件等价：
\end{enumerate}

\begin{enumerate}
\def\labelenumi{\alph{enumi})}
\item
  A 是正则的；
\item
  A 是整闭的；
\item
  A 是离散赋值环。
\end{enumerate}

\begin{enumerate}
\def\labelenumi{\arabic{enumi})}
\setcounter{enumi}{2}
\tightlist
\item
  若对于 A 的每个素理想 p，环 \(\mathbf{A}_{\mathfrak{p}}\) 是整闭的，则称环 A 是正规的。
\end{enumerate}

\begin{enumerate}
\def\labelenumi{\alph{enumi})}
\item
  正规环 A 是约化的，并且每个素理想 p 恰好包含 A 的一个极小素理想。对于每个极小素理想 p，环 A/p 是整闭的。
\item
  诺特正规环 A 同构于有限个整闭诺特环的直积。
\item
  设 A 是诺特正规环，a 是 A 的一个元素。则 Ass \(_{A}\) (A/aA) 不含嵌入素理想（IV，§ 2，第 3 号，注）。
\end{enumerate}

\begin{enumerate}
\def\labelenumi{\arabic{enumi})}
\setcounter{enumi}{3}
\tightlist
\item
  设 A 是一个环。我们考察 A 上的以下两个条件：
\end{enumerate}

(R1) 对 A 的每个高度 ≤ 1 的素理想 \(\mathfrak{p}\)，环 \(A_{\mathfrak{p}}\) 是正则的。

(S2) 集合 \(\mathrm{Ass}_{\mathbf{A}}(\mathbf{A})\)，以及对于 A 的每个可消去元素 \(a\) 的集合 \(\mathrm{Ass}_{\mathbf{A}}(\mathbf{A}/a\mathbf{A})\)，不含嵌入素理想。

  诺特正规环（习题 3）满足 (R1) 和 (S2)。（使用习题 2 和 3(c)。）

\begin{enumerate}
\def\labelenumi{\arabic{enumi})}
\setcounter{enumi}{4}
\tightlist
\item
  设 A 是满足条件 (R1) 和 (S2) 的诺特环。
\end{enumerate}

\begin{enumerate}
\def\labelenumi{\alph{enumi})}
\item
  证明 A 是约化的。
\item
  设 \(a\) 是 A 的一个可消去元素。如果 A 的一个元素 \(x\) 的像在 \(\mathbf{A}_{\mathfrak{p}}\) 中属于 \(a\mathbf{A}_{\mathfrak{p}}\)，且这对于每个 \(\mathfrak{p} \in \operatorname{Ass}_{\mathbf{A}}(\mathbf{A}/a\mathbf{A})\) 都成立，则 \(x\) 属于 \(a\mathbf{A}\)。（使用 IV，§ 2，第 3 号，命题 5。）
\item
  证明 A 在其总分式环 R 中是整闭的。（若 \(a, b, c_{1}, \ldots, c_{n}\) 是 A 中满足以下条件的元素
\end{enumerate}

\begin{enumerate}
\def\labelenumi{(\roman{enumi})}
\item
  \(b\) 在 A 中是可消去的，
\item
  \((b^{-1}a)^{n} + c_{1}(b^{-1}a)^{n-1} + \cdots + c_{n} = 0\) 在 R 中，则依次证明有 \(a \in bA_{\mathfrak{p}}\)，其中 \(\mathfrak{p}\) 遍历 A 中的每个高度为 1 的素理想，进而 \(a \in bA\)。）
\end{enumerate}

\begin{enumerate}
\def\labelenumi{\alph{enumi})}
\setcounter{enumi}{3}
\tightlist
\item
  证明 A 是正规的。（注意 R 中的幂等元在 A 上是整的。）
\end{enumerate}

\begin{enumerate}
\def\labelenumi{\arabic{enumi})}
\setcounter{enumi}{5}
\tightlist
\item
  \begin{enumerate}
  \def\labelenumii{\alph{enumii})}
  \tightlist
  \item
    设 A 是一个环，M 是忠实的诺特 A-模。证明 A 是诺特环。
  \end{enumerate}
\end{enumerate}

\begin{enumerate}
\def\labelenumi{\alph{enumi})}
\setcounter{enumi}{1}
\item
  设 A 是一个环，M 是忠实的有限生成 A-模。假设对于 A 的每个理想的递增序列 \((\mathfrak{a}_{i})_{i\in\mathbb{N}}\)，M 的子模序列 \((\mathfrak{a}_{i}\mathbf{M})_{i\in\mathbb{N}}\) 是稳定的。则 A 是诺特环。（根据 (a)，只需证明 M 是诺特 A-模。将情形化归到：对于 A 的每个非零理想 a，A-模 M/aM 是诺特的；并且对于 M 的每个非零子模 N，A-模 M/N 不是忠实的。）
\item
  设 B 是一个诺特环，A 是 B 的子环，且 B 是有限 A-代数。则 A 是诺特环。（这使得可以在 § 4 第 1 号的命题 2 中去掉 A 是诺特环的假设。）
\end{enumerate}

\begin{enumerate}
\def\labelenumi{\arabic{enumi})}
\setcounter{enumi}{6}
\tightlist
\item
  设 A 是一个 Krull 环，P 是其高度 1 素理想的集合。假设对于每个素理想 \(p \in P\)，环 A/p 是诺特的，并记 K 为 A 的分式域。证明 A 是诺特的。
\end{enumerate}

(设 \(\mathfrak{p} \in \mathrm{P}\)，且 \(x \in \mathbf{K}\) 满足 \(v_{\mathfrak{p}}(x) = 1\) 及 \(v_{\mathfrak{q}}(x) \leqslant 0\)，其中 \(\mathfrak{q} \in \mathrm{P} - \{\mathfrak{p}\}\)。置 \(\mathbf{B} = \mathbf{A}[x]\)。则有以下性质：

\begin{enumerate}
\def\labelenumi{(\roman{enumi})}
\item
  \(\mathfrak{p} = x\mathbf{B}\cap \mathbf{A};\)
\item
  A 到 B 的包含映射通过取商诱导出从 A/p 到 B/xB 的同构；
\item
  对于每个整数 \(n \geqslant 0\)，环 \(\mathbf{B}/x^{n}\mathbf{B}\) 是诺特的；
\item
  对于每个整数 \(n \geqslant 0\)，环 \(\mathbf{A}/(x^n\mathbf{B} \cap \mathbf{A})\) 是诺特的（使用习题 6）；
\item
  对于每个整数 \(n \geqslant 0\)，A-模 \(\mathbf{A}/\mathfrak{p}^{(n)}\) 是诺特的。回忆我们定义 \(\mathfrak{p}^{(n)} = \mathbf{A} \cap \mathfrak{p}^n\mathbf{A}_{\mathfrak{p}}\)（参见 IV，§ 2，习题 18。）
\end{enumerate}

\begin{enumerate}
\def\labelenumi{\arabic{enumi})}
\setcounter{enumi}{7}
\tightlist
\item
  设 A 是一个整环，K 是其分式域。设 B 是一个环，且 \(\varphi: A \to B\) 是使 B 成为忠实平坦 A-模的环同态。假设 B 只有有限多个极小素理想 \(\mathfrak{p}_1, \ldots, \mathfrak{p}_n\)。对于 \(1 \leqslant i \leqslant n\)，置 \(\mathbf{B}_i = \mathbf{B} / \mathfrak{p}_i\)，记 \(\mathbf{L}_i\) 为 \(\mathbf{B}_i\) 的分式域，并记 L 为各 \(\mathbf{L}_i\) 的直积环。最后，记 \(\overline{\mathbf{A}}\) 为 A 的整闭包，\(\overline{\mathbf{B}}_i\) 为 \(\mathbf{B}_i\) 的整闭包。
\end{enumerate}

\begin{enumerate}
\def\labelenumi{\alph{enumi})}
\item
  由 \(\varphi\) 导出的从 A 到 L 的同态是单射，并且延拓为从 K 到 L 的单射同态 \(\psi\)。
\item
  有 \(\overline{\mathbf{A}} = \psi^{-1}(\prod_{i=1}^{n} \overline{\mathbf{B}}_i)\)。（可以证明：若 \(x \in \psi^{-1}(\prod_{i=1}^{n} \overline{\mathbf{B}}_i)\)，且 \(x'\) 是它在 \(\mathbf{K} \otimes_{\mathbf{A}} \mathbf{B}\) 中的像，则存在一个首一多项式 \(\mathbf{P}\)，其系数在 \(\mathbf{B}\) 中，使得 \(\mathbf{P}(x')\) 在 \(\mathbf{K} \otimes_{\mathbf{A}} \mathbf{B}\) 中是幂零的。）
\end{enumerate}

\begin{enumerate}
\def\labelenumi{\arabic{enumi})}
\setcounter{enumi}{8}
\tightlist
\item
  设 A 是一个诺特整环，K 是其分式域，E 是 K 的有限扩张，且 \(\overline{\mathbf{A}}\) 是 A 在 E 中的整闭包。
\end{enumerate}

\begin{enumerate}
\def\labelenumi{\alph{enumi})}
\item
  若 A 是带有完备化 \(\hat{\mathbf{A}}\) 的局部环，则 \((\hat{\mathbf{A}} \otimes_{\mathbf{A}} \overline{\mathbf{A}})_{\mathrm{red}}\) 是有限 \(\hat{\mathbf{A}}\)-代数 \(^{1}\)（将情形化归到 \(\mathbf{K} = \mathbf{E}\) 的情形。利用 III，§ 3，第 4 号，定理 3 的推论 2，证明 A 的非零元素在 \(\hat{\mathbf{A}}_{\mathrm{red}}\) 中都不是零因子。然后证明存在一个 \(\hat{\mathbf{A}}_{\mathrm{red}}\)-单射同态，从 \((\hat{\mathbf{A}} \otimes_{\mathbf{A}} \overline{\mathbf{A}})_{\mathrm{red}}\) 到 \(\hat{\mathbf{A}}_{\mathrm{red}}\) 的总分式环中。结论。）
\item
  对于 A 的每个素理想 p，\(\kappa(p)\)-代数 \((\overline{A} \otimes_A \kappa(p))_{\text{red}}\) 是有限的。（将情形化归到 A 是以 p 为极大理想的局部环的情形。然后注意到 \((\overline{A} \otimes_A \kappa(p))_{\text{red}}\) 是 \((\overline{A} \otimes_A \hat{A})_{\text{red}} \otimes \kappa(p).)\) 的一个商。）
\end{enumerate}

\begin{enumerate}
\def\labelenumi{\arabic{enumi})}
\setcounter{enumi}{9}
\tightlist
\item
  设 A 是一个局部诺特整环，\(\hat{\mathbf{A}}\) 是其完备化，\(\mathfrak{p}_1, \ldots, \mathfrak{p}_n\) 是 \(\hat{\mathbf{A}}\) 的极小素理想，并且对于 \(1 \leqslant i \leqslant n\)，置 \(\mathbf{B}_i = \hat{\mathbf{A}} / \mathfrak{p}_i\)。
\end{enumerate}

\begin{enumerate}
\def\labelenumi{\alph{enumi})}
\item
  环 \(\mathbf{B}_i\) 是日式的。
\item
  \(\overline{\mathbf{B}}_{i}\) 是 \(\mathbf{B}_{i}\) 的整闭包，且是 Krull 环。
\item
  \(\overline{\mathbf{A}}\) 是 \(\mathbf{A}\) 的整闭包，且是 Krull 环。（使用习题 8 和 VII，§ 1，第 3 号，例 3 和 4。）
\item
  \(\overline{\mathrm{A}}\) 只有有限个极大理想，并且它们的剩余域是 A 的剩余域的有限扩张。（使用习题 9(b)。）
\item
  对于 A 的每个素理想 p，\(\overline{A}\) 在 p 之上的素理想 P 只有有限多个，并且域 \(\kappa(\mathfrak{P})\) 是 \(\kappa(p)\) 的有限扩张（将情形化归到 A 是以 p 为极大理想的局部环的情形。）
\end{enumerate}

\begin{enumerate}
\def\labelenumi{\arabic{enumi})}
\setcounter{enumi}{10}
\tightlist
\item
  设 A 是一个诺特整环，K 是其分式域，E 是 K 的有限扩张，\(\overline{A}\) 是 A 在 E 中的整闭包。
\end{enumerate}

\begin{enumerate}
\def\labelenumi{\alph{enumi})}
\item
  对于 A 的每个素理想 p，\(\overline{A}\) 在 p 之上的素理想 P 只有有限多个，且域 \(\kappa(\mathfrak{P})\) 是 \(\kappa(p)\) 的有限扩张。（将情形化归到 A 是以 p 为极大理想且 E = K 的情形，并使用习题 10。）
\item
  设 P 是 \(\overline{\mathbf{A}}\) 的高度为 1 的素理想集合。如果 \(p \in P\)，则 \(\overline{\mathbf{A}}_p\) 是离散赋值环。
\item
  有 \(\overline{\mathbf{A}} = \bigcap_{\mathfrak{p} \in \mathbb{P}} \overline{\mathbf{A}}_{\mathfrak{p}}\)（在 E 中）。
\end{enumerate}

\begin{enumerate}
\def\labelenumi{\arabic{enumi})}
\setcounter{enumi}{11}
\tightlist
\item
  设 A 是一个整环，\(\overline{\mathbf{A}}\) 是其整闭包。设 \(a_1, \ldots, a_r\) 是 \(\overline{\mathbf{A}}\) 的非零元素，且 b 是 A 的分式理想 \(\sum_{i=1}^{r} \mathbf{A}a_i\)。
\end{enumerate}

\begin{enumerate}
\def\labelenumi{\alph{enumi})}
\item
  存在 A 的一个分式理想 \(\mathfrak{a}\)，使得 \(\mathfrak{a}a_{i} \subset \mathfrak{a}\) 对于 \(1 \leqslant i \leqslant r\) 成立。我们还有 \(\tilde{\mathfrak{a}}a_{i} \subset \tilde{\mathfrak{a}}\) 对于 \(1 \leqslant i \leqslant r\) 成立；回忆（VII，§ 1，第 1 号），\(\tilde{\mathfrak{a}}\) 是包含 \(\mathfrak{a}\) 的分式主理想的交。
\item
  设 y 是 \(\tilde{\mathfrak{b}}\) 的一个非零元素。则
\end{enumerate}

\[
\mathrm{A} y ^ {- 1} \supset \bigcap_ {i = 1} ^ {r} \mathrm{A} a _ {i} ^ {- 1} \quad \text { and } \quad \widetilde {\mathfrak {a}} \subset \widetilde {\mathfrak {a}} y ^ {- 1}.
\]

\begin{enumerate}
\def\labelenumi{\alph{enumi})}
\setcounter{enumi}{2}
\item
  有 \(\tilde{\mathbf{b}}\subset \overline{\mathbf{A}}\)
\item
  设 z 是 \(\overline{\mathbf{A}}:\overline{\mathbf{A}}\mathbf{b}\) 的一个元素。则 \(\overline{\mathbf{A}}\) 包含 \(\tilde{\mathbf{b}}z\)。
\item
  有 b ⊂ \(\overline{A}:(\overline{A}:\overline{A}b)\)。
\item
  设 \(\mathfrak{P}\) 是 \(\overline{\mathbf{A}}\) 的一个除子理想，且 \(\mathfrak{p} = \mathfrak{P} \cap \mathbf{A}\)。则 \(\mathfrak{p}\) 是 \(\mathbf{A}\) 中的除子理想。（由 \(e\)），得到 \(\tilde{\mathfrak{p}} \subset \mathfrak{P}\)。但我们还有 \(\tilde{\mathfrak{p}} \subset \mathbf{A}\)。）
\end{enumerate}

\begin{enumerate}
\def\labelenumi{\arabic{enumi})}
\setcounter{enumi}{12}
\tightlist
\item
  设 A 是一个诺特整环，K 是其分式域，\(\overline{\mathbf{A}}\) 是其整闭包，\(f\) 是 A 的一个非零元素，\(\mathfrak{P}\) 是 \(\overline{\mathbf{A}}\) 的一个包含 \(f\) 的高度为 1 的素理想，并且设 \(\mathfrak{p} = \mathbf{A} \cap \mathfrak{P}\) 是这个素理想
\end{enumerate}

\(^{1}\) 对于每个环 B，我们将 \(B_{red}\) 记为 B 关于其幂零元理想的商。

的素理想。证明 \(\mathfrak{p} \in \operatorname{Ass}_{\mathbf{A}}(\mathbf{A}/f\mathbf{A})\)，方法是将情形化归到 A 是以 \(\mathfrak{p}\) 为极大理想的局部环，并在此假设下依次建立以下断言：a) \(\mathfrak{P}\) 是除子理想（使用习题 10(c)）。

\begin{enumerate}
\def\labelenumi{\alph{enumi})}
\setcounter{enumi}{1}
\item
  p 是除子理想（使用习题 12(f)）。
\item
  若 \(a_{1},\ldots,a_{r}\) 是 K 的非零元素，使得
\end{enumerate}

\[
\mathrm{A}: \mathfrak {p} = \mathrm{A} + \sum_ {i = 1} ^ {r} \mathrm{A} a _ {i}, \quad \text { and } \quad \mathfrak {p} = \bigcap_ {i = 1} ^ {r} (\mathrm{A} \cap a _ {i} ^ {- 1} \mathrm{A})
\]

并且存在一个指标 \(i\)，使得 \(\mathfrak{p} = \mathbf{A} \cap a_i^{-1}\mathbf{A}\)。

\begin{enumerate}
\def\labelenumi{\alph{enumi})}
\setcounter{enumi}{3}
\item
  存在 A 的一个非零元素 g，使得 \(\mathfrak{p} \in \operatorname{Ass}_{\mathbf{A}}(\mathbf{A}/g\mathbf{A})\)。
\item
  \(\mathfrak{p} \in \operatorname{Ass}_{\mathbf{A}}(\mathbf{A}/f\mathbf{A})\)。
\end{enumerate}

\begin{enumerate}
\def\labelenumi{\arabic{enumi})}
\setcounter{enumi}{13}
\item
  设 A 是一个诺特整环，K 是其分式域，\(\overline{\mathbf{A}}\) 是其在 K 的有限扩张 E 中的整闭包。则 \(\overline{\mathbf{A}}\) 是 Krull 环。（使用习题 11(c)；并且为了证明每个元素 \(f\)（属于 \(\overline{\mathbf{A}}\) 且非零）只属于有限多个高度为 1 的素理想，将情形化归到 \(f \in \mathbf{A}\)、\(\mathbf{E} = \mathbf{K}\) 的情形，并使用习题 13(b) 和 11(a)。）
\item
  设 A 是一个诺特环，B 是一个整 A-代数，只有有限多个极小素理想 \(p_{1}, \ldots, p_{n}\)。对于 \(1 \leqslant i \leqslant n\)，将 \(q_{i}\) 定义为 \(p_{i}\) 在 A 中的逆像，并假设 \(\kappa(p_{i})\) 是 \(\kappa(q_{i})\) 的有限扩张。则对于 B 的每个元素 x，Spec(B) 的拓扑子空间 V(x) 只有有限多个不可约分支。（将情形化归到 B 是整闭且包含 A 的情形，并使用习题 14。）
\item
  设 A 是一个诺特整环，\(\overline{\mathbf{A}}\) 是 A 在其分式域的一个有限扩张中的整闭包，B 是一个包含 A 且包含于 \(\overline{\mathbf{A}}\) 的环。置 \(\mathbf{Y} = \operatorname{Spec}(\mathbf{B})\) 和 \(\mathbf{X} = \operatorname{Spec}(\mathbf{A})\)。我们拟证明 Y 是一个诺特空间。设 \((f_n)_{n \in \mathbb{N}}\) 是 B 的元素序列。对于 \(n \in \mathbb{N}\)，令 \(\mathbf{U}_n = \operatorname{Spec}(\mathbf{B}_{f_n})\) 为由 \(f_n\) 定义的 Y 的开子集。置 \(\mathbf{U} = \bigcup_{n \in \mathbb{N}} \mathbf{U}_n\)，\(\mathbf{F}_n = \mathbf{U} - \bigcup_{0 \leqslant m \leqslant n} \mathbf{U}_m\)。我们把 \(\overline{\mathbf{F}}_n\) 记为 \(\mathbf{F}_n\) 在 Y 中的闭包，把 \(\mathbf{G}_n\) 记为 \(\mathbf{F}_n\)
\end{enumerate}

在 X 中的像，把 \(\overline{\mathbf{G}}_n\) 记为 \(\mathbf{G}_n\) 在 X 中的闭包，把 \(\mathbf{B}_n\) 记为满足 \(\operatorname{Spec}(\mathbf{B}_n) = \overline{\mathbf{F}}_n\) 的 B 的约化商环，把 \(\overline{f}_n\) 记为 \(f_n\) 在 \(\mathbf{B}_n\) 中的像。

\begin{enumerate}
\def\labelenumi{\alph{enumi})}
\item
  有 \(\mathbf{F}_n = \mathbf{V}(\overline{f}_n\mathbf{B}_n)\cap \mathbf{U}\) 对于每个 \(n\in \mathbf{N}\)。
\item
  假设对于某个整数 \(n \in N\)，\(\overline{F}_{n}\) 只有有限多个不可约分支。则 \(V(\overline{f}_{n}B_{n})\) 只有有限多个不可约分支。（应用习题 15 的结果，同时使用习题 11(a)。）
\item
  对于每个 \(n\in\mathbf{N}\)，\(\mathbf{F}_{n}\) 只有有限多个不可约分支。
\item
  若 \(\mathbf{A}_n\) 是满足 \(\operatorname{Spec}(\mathbf{A}_n) = \overline{\mathbf{G}}_n\) 的 A 的约化商环，则 \(\mathbf{A}_n\) 的极小素理想属于 \(\mathbf{G}_n\)。
\item
  令 \(n \in \mathbf{N}\)，且 \(\mathfrak{x}\) 是 \(\mathbf{A}_n\) 的这样一个极小素理想。则存在一个整数 \(n' > n\)，使得 \(\mathbf{F}_{n'}\) 不包含 \(\mathbf{Y}\) 中位于 \(\mathfrak{x}\) 之上的点，并且 \(\mathfrak{x}\) 不属于 \(\overline{\mathbf{G}}_{n'}\)。（使用习题 11(a)。）
\item
  当 n 足够大时，\(F_{n}\) 为空。
\end{enumerate}

\begin{enumerate}
\def\labelenumi{\arabic{enumi})}
\setcounter{enumi}{16}
\tightlist
\item
  设 A 是一个诺特整环，且 \(a \neq 0\) 是 A 的根基中的元素。我们作如下假设：
\end{enumerate}

\begin{enumerate}
\def\labelenumi{(\roman{enumi})}
\item
  \(\operatorname{Ass}_{\mathbf{A}}(\mathbf{A} / a\mathbf{A})\) 包含唯一一个极小元素 \(p\)；
\item
  \(a\mathbf{A}_{\mathfrak{p}} = \mathfrak{p}\mathbf{A}_{\mathfrak{p}}\)
\item
  A/p 是整闭的；
\item
  A 的整闭包 \(\overline{\mathbf{A}}\) 是有限 A-代数；
\item
  若 \(\mathfrak{q}\) 是 \(\overline{\mathbf{A}}\) 中高度为 1 的素理想，则 \(\mathfrak{q} \cap A\) 是 A 中高度为 1 的素理想。
\end{enumerate}

证明 \(\overline{\mathbf{A}} = \mathbf{A}\)（A 是整闭的）以及 \(a\mathbf{A} = \mathfrak{p}\)（因而 A/aA 是整闭的）。为此依次建立以下断言：

\begin{enumerate}
\def\labelenumi{\alph{enumi})}
\item
  \(A_{p}\) 与 \(\overline{A}_{p}\) 同构，且 \(\overline{p} = p\overline{A}\) 是 \(\overline{A}\) 位于 p 之上的唯一素理想。
\item
  此外，\(\overline{p}\) 是 \(\operatorname{Ass}_{\mathbf{A}}(\overline{\mathbf{A}}/a\overline{\mathbf{A}})\) 的唯一极小元素。
\item
  事实上，\(\overline{p}\) 是 \(\operatorname{Ass}_{\mathbf{A}}(\overline{\mathbf{A}}/a\overline{\mathbf{A}})\) 的唯一元素。
\item
  \(\overline{\mathbf{A}} / a\overline{\mathbf{A}}\) 是整环，且 \(a\overline{\mathbf{A}} = \overline{\mathfrak{p}}\)。
\item
  A/p 与 \(\overline{A}/\overline{p}\) 同构，且从 A/aA 到 \(\overline{A}/a\overline{A}\) 的典范同态是满射的。
\end{enumerate}

\begin{enumerate}
\def\labelenumi{\arabic{enumi})}
\setcounter{enumi}{17}
\tightlist
\item
  设 A 是一个诺特整环，\(x\) 是 A 的非零元素。我们假设 A 关于 \(xA\)-进拓扑是分离且完备的，并且对于每个素理想 \(p \in \mathrm{Ass}_{\mathbf{A}}(\mathbf{A} / x\mathbf{A})\)，环 \(A/p\) 是日式的。证明 A 是日式的。（设 \(\overline{\mathbf{A}}\) 是 A 在其分式域的一个有限扩张中的整闭包。
\end{enumerate}

\begin{enumerate}
\def\labelenumi{\alph{enumi})}
\item
  由于 \(\overline{\mathbf{A}}\) 是 Krull 环（第 81 页，习题 14），设 \(\mathfrak{P}_1, \ldots, \mathfrak{P}_n\) 是 \(\overline{\mathbf{A}}\) 中包含 \(x\) 的高度为 1 的素理想。
\item
  令 \(\mathfrak{p}_i = \mathfrak{P}_i \cap \mathbf{A}\)。则有 \(\mathfrak{p}_i \in \operatorname{Ass}_{\mathbf{A}}(\mathbf{A} / x\mathbf{A})\)，且 \(\kappa(\mathfrak{P}_i)\) 是 \(\kappa(\mathfrak{p}_i)\) 的有限扩张（使用第 80 页习题 11 和 13。）
\item
  \(\overline{\mathbf{A}} / \mathfrak{P}_i\) 是有限生成的 \(\mathbf{A} / \mathfrak{p}_i\)-模。
\item
  \(\overline{\mathbf{A}} / x\overline{\mathbf{A}}\) 是有限生成的 A/xA-模。
\item
  \(\overline{\mathbf{A}}\) 关于 \(x\overline{\mathbf{A}}\)-进拓扑是分离的。
\end{enumerate}

如 § 4 第 2 号定理 1 的证明那样作结。

\begin{enumerate}
\def\labelenumi{\arabic{enumi})}
\setcounter{enumi}{18}
\tightlist
\item
  设 A 是一个诺特环，a 是 A 的非零理想。假设 A 关于 a-进拓扑是分离且完备的，且 A/a 是 Nagata 环。证明 A 是 Nagata 环。可以如下论证：
\end{enumerate}

\begin{enumerate}
\def\labelenumi{\alph{enumi})}
\item
  假设 A 是一个整环，并且对于 A 的每个非零素理想 p，环 A/p 是日式的。设 \(\overline{\mathbf{A}}\) 是 A 在其分式域的一个有限扩张中的整闭包。则对于 \(\mathfrak{P}\)（\(\overline{\mathbf{A}}\) 的每个非零素理想），环 \(\overline{\mathbf{A}}/\mathfrak{P}\) 是环 \(A/(\mathfrak{P} \cap A)\) 上的有限代数。（使用第 80 页习题 11(a)。）
\item
  在与 a) 中相同的假设和记号下，证明 \(\overline{A}\) 是诺特的。（使用第 79 页习题 7 和第 81 页习题 14。）然后仿照前一个习题的推理，证明 \(\overline{A}/a\overline{A}\) 是有限生成的 A-模，\(\overline{A}\) 关于 \(a\overline{A}\) 拓扑是分离的，并且 \(\overline{A}\) 是有限生成的 A-模。
\end{enumerate}

\begin{enumerate}
\def\labelenumi{\arabic{enumi})}
\setcounter{enumi}{19}
\item
  设 A 是一个诺特环，a 是 A 的一个不同于 A 的理想，\(\hat{A}\) 是 A 关于 a-进拓扑的分离完备化。如果 A/a 是 Nagata 环，则 \(\hat{A}\) 是 Nagata 环。（使用习题 19。）
\item
  \begin{enumerate}
  \def\labelenumii{\alph{enumii})}
  \tightlist
  \item
    设 A 是一个整环，K 是其分式域。如果 K 的特征 p 非零，并且 E 是有理函数域 \(\mathrm{K}(\mathrm{X})\) 的有限纯不可分扩张，则存在 K 的有限纯不可分扩张 F 和 p 的一个幂 q，使得 \(\mathrm{F}(\mathrm{X}^{1/q})\) 是 E 的一个扩张。
  \end{enumerate}
\end{enumerate}

\begin{enumerate}
\def\labelenumi{\alph{enumi})}
\setcounter{enumi}{1}
\item
  设 A 是整闭的 Nagata 环。则多项式环 A{[}X{]} 是日式的。
\item
  设 A 是整闭的 Nagata 环。则有限个不定元上的每个多项式环 \(A[X_1, \ldots, X_n]\) 都是 Nagata 环。
\end{enumerate}

\begin{enumerate}
\def\labelenumi{\arabic{enumi})}
\setcounter{enumi}{21}
\item
  设 A 是一个诺特环，a 是 A 的一个不同于 A 的理想。假设 A 关于 a-进拓扑是分离且完备的，且 A/a 是 Nagata 环。则有限个不定元上的每个限制形式幂级数环都是 Nagata 环。（使用 III，§ 4，习题 7，以及上面的习题 20 和 21。）
\item
  设 A 是约化诺特半局部环，R 是其总分式环，\(\hat{A}\) 是 A 的完备化。要使 \(\hat{A}\) 约化，充要条件是 \(R \otimes_{A} \hat{A}\) 是约化的。
\item
  设 A 是半局部诺特环，\(\hat{\mathbf{A}}\) 是其完备化，\(x\) 是 A 的根基中的非零元素。我们假设 \(\mathrm{Ass}_{\mathbf{A}}(\mathbf{A}/x\mathbf{A})\) 不含嵌入素理想，并且对于每个 \(\mathfrak{p} \in \mathrm{Ass}_{\mathbf{A}}(\mathbf{A}/x\mathbf{A})\)，环 \(\mathbf{A}_{\mathfrak{p}}\) 是正则的，且环 \(\kappa(\mathfrak{p}) \otimes_{\mathbf{A}} \hat{\mathbf{A}}\) 是约化的。则 \(\hat{\mathbf{A}}\) 是约化的。（理由同 § 4 第 4 号定理 3 中 (i) \(\Rightarrow\) (ii) 的证明。）
\item
  设 A 是一个诺特整环，X 是拓扑空间 \(\operatorname{Spec}(A)\)，Nor(X) 是 X 中的点 \(\mathfrak{p}\) 的集合，这些点满足局部环 \(A_{\mathfrak{p}}\) 是整闭的。（如果 A 的整闭包是有限 A-代数，则 Nor(X) 在 X 中是开的，参见 V，§ 1，第 5 号，推论 5。）
\end{enumerate}

令 \(f\) 是 A 的一个非零元素，使得环 \(\mathbf{A}[f^{-1}]\) 是整闭的。

\begin{enumerate}
\def\labelenumi{\alph{enumi})}
\item
  若 A 的一个素理想 p 不包含 f，则 p 属于 Nor(X)。
\item
  设 E 是 A 的素理想 \(\mathfrak{p}\) 的集合，这些素理想与 A/fA 相关联，其中或者高度为 \textgreater{} 1，或者高度为 1 且 \(A_{\mathfrak{p}}\) 不是正则的。则有
\end{enumerate}

\[
\operatorname{Nor} (X) = X - \bigcup_ {\mathfrak {p} \in E} V (\mathfrak {p}).
\]

（使用第 79 页习题 4 和 5。）

\begin{enumerate}
\def\labelenumi{\alph{enumi})}
\setcounter{enumi}{2}
\tightlist
\item
  Nor(X) 在 X 中是开的。
\end{enumerate}

\begin{enumerate}
\def\labelenumi{\arabic{enumi})}
\setcounter{enumi}{25}
\tightlist
\item
  设 A 是一个诺特整环，X 是拓扑空间 \(\operatorname{Spec}(A)\)，Nor(X) 是习题 25 中引入的 X 的子空间，\(\overline{A}\) 是 A 的整闭包。我们假设存在 A 的一个非零元素 \(f\)，使得环 \(A[f^{-1}]\) 是整闭的，并且对于 A 的每个极大理想 m，环 \(\overline{A}_{m}\) 是有限 A \(_{m}\)-代数。把 \(\overline{A}\) 看作有限子 A-代数的滤增归纳极限，\(\overline{A} = \varinjlim(A_{j})_{j\in J}\)。置 \(X_{j} = \operatorname{Spec}(A_{j})\)，并令 \(G_{j}\) 为 \(X_{j} - \operatorname{Nor}(X_{j})\) 在 X 中的像。
\end{enumerate}

\begin{enumerate}
\def\labelenumi{\alph{enumi})}
\item
  \(\operatorname{Nor}(\mathbf{X}_j)\) 在 \(\mathbf{X}_j\) 中是开的。（使用习题 25。）
\item
  \(G_{j}\) 在 X 中是闭的。
\item
  存在一个指标 \(j_0 \in J\)，使得 \(G_{j_0} = \bigcap G_j\)。
\end{enumerate}

\(d)\) 令 \(\mathfrak{p} \in \operatorname{Spec}(\mathbf{A})\)。则对于充分大的 \(j \in \mathrm{J}\)，\(\mathbf{G}_j\) 不包含 \(\mathfrak{p}\)。

\begin{enumerate}
\def\labelenumi{\alph{enumi})}
\setcounter{enumi}{4}
\tightlist
\item
  \(\overline{\mathbf{A}}\) 是有限 A-代数。
\end{enumerate}

\begin{enumerate}
\def\labelenumi{\arabic{enumi})}
\setcounter{enumi}{26}
\tightlist
\item
  设 A 是诺特整闭局部环。此外假设 A 是 Nagata 环。设 x 是 A 的分式域 K 中不属于 A 的元素。设 B 是环 A{[}x{]}，且 p 是包含 x 的 B 的极大理想。置 I = xA ∩ A。
\end{enumerate}

\begin{enumerate}
\def\labelenumi{\alph{enumi})}
\item
  设 X 是一个不定元，令 Q 为将 X 映到 x 的 A{[}X{]} 到 B 的同态的核。则 Q 由形如 aX - b 的多项式生成，其中 a 和 b 是 A 中满足 ax = b 的元素，并且 I = xA ∩ A 由这些多项式的常数项 b 生成。此外，A 到 B 的包含映射诱导出从 A/I 到 B/xB 的同构。
\item
  A 的理想 I 是除子理想。
\item
  \(\operatorname{Ass}_{\mathbf{B}}(\mathbf{B}/x\mathbf{B})\) 不含嵌入素理想。
\item
  设 q 是 \(B_{p}\) 的素理想，且与 \(B_{p}/xB_{p}\) 相关联。则 \(q \cap A\) 与 A/I 相关联，且环 \(A_{q \cap A}\) 和 \((\mathbf{B}_{\mathfrak{p}})_{q}\) 是离散赋值环。
\item
  \(B_{p}/q\) 是 Nagata 环。（注意 \(\mathrm{B}/(\mathrm{q} \cap \mathrm{B})\) 同构于 \(\mathrm{A}/(\mathrm{q} \cap \mathrm{A})\)。）
\item
  \(B_{p}/q\) 的完备化是约化的。
\item
  \(B_{p}\) 的完备化是约化的，且 \(B_{p}\) 的整闭包是有限 \(B_{p}\)-代数。（使用习题 24。）
\end{enumerate}

\begin{enumerate}
\def\labelenumi{\arabic{enumi})}
\setcounter{enumi}{27}
\tightlist
\item
  设 A 是诺特整闭局部环，K 是其分式域，x 是 K - A 中的元素，B 是环 A{[}x{]}，且 p 是 B 的一个极大理想。设 a、b 是 A 的非零元素，满足 bx = a。
\end{enumerate}

\begin{enumerate}
\def\labelenumi{\alph{enumi})}
\item
  环 B \(\left[\frac{1}{b}\right]\) 是整闭的。
\item
  存在一个系数在 A 中的首一多项式 \(f\)，使得 \(f(x)\in\mathfrak{p}\)。
\item
  设 E 是通过将 \(f(X)\) 的根添加到 K 所得的域，A' 是 A 在 E 中的整闭包，B' 是环 A'{[}x{]}。设 \(\mathfrak{p}'\) 是位于 p 之上的 B' 的一个极大理想。则 \(B'_{\mathfrak{p}'}\) 的整闭包是有限 \(B'_{\mathfrak{p}'}\)-代数。（使用前一个习题。）
\item
  \(\mathbf{B}_{\mathfrak{p}'}^{\prime}\) 的整闭包是有限 \(\mathbf{B}_{\mathfrak{p}'}^{\prime}\)-代数。（使用习题 26。）
\item
  \(B_{p}\) 的整闭包是有限 \(B_{p}\)-代数。
\end{enumerate}

\begin{enumerate}
\def\labelenumi{\arabic{enumi})}
\setcounter{enumi}{28}
\tightlist
\item
  设 A 是整闭的 Nagata 环，x 是 A 的分式域中不属于 A 的元素。设 B 是环 A{[}x{]}，且 a、b 是 A 的两个非零元素，满足 bx = a。
\end{enumerate}

\begin{enumerate}
\def\labelenumi{\alph{enumi})}
\item
  环 \(\mathbf{B}\left[\frac{1}{b}\right]\) 是整闭的。
\item
  对于 B 的每个极大理想 m，\(B_{m}\) 的整闭包是有限 \(B_{m}\)-代数。（使用前一个习题。）
\item
  B 的整闭包是有限 B-代数。（使用习题 26。）
\end{enumerate}

\begin{enumerate}
\def\labelenumi{\arabic{enumi})}
\setcounter{enumi}{29}
\item
  设 A 是 Nagata 环。则 A 上的每个有限型代数都是 Nagata 环。（使用第 82 页习题 21 和 29。）
\item
  设 A 是一个诺特整环，\(\overline{A}\) 是 A 在其分式域的一个有限扩张中的整闭包。如果 A 的维数至多为 2，则 \(\overline{A}\) 是诺特环。（使用第 79 页习题 7、第 81 页习题 14 以及 Krull-Akizuki 定理（参见 VII，§ 2，第 5 号）。）
\item
  设 A 是局部诺特整环，K 是其分式域。假设 K 的特征 p 非零。假设 A 是 Nagata 环，并记 \(\hat{A}\) 为 A 的完备化。设 p 是 \(\hat{A}\) 的一个极小素理想，L 是 \(\hat{A}/p\) 的分式域。令 k 是 A 的弱剩余域，\(k'\) 是 \(\hat{A}\) 的代表元域，包含 k 在 \(\hat{A}\) 中的像（第 78 页，习题 9(a)。）。则 K 是 k 的可分扩张，当且仅当 L 是 \(k'\) 的可分扩张。（注意 L 是 K 的可分扩张，因为 A 是 Nagata 环。如果 K 是 k 的可分扩张，证明从 \(k'\) 到 L 的每个导数都延拓为从 L 到 L 的导数。）
\end{enumerate}

